% SOURCE AUTHORITY:
% expose_VII_body.tex lines 941--984; scan indices 185--187; source folios 174--176.
% Source-backed corrections: references to "2.9.2 b)" are to Proposition 2.9.3(b);
% Biext in (2.9.6.1) is written Biext^1 to identify the group of isomorphism classes.
% Hard stop before source line 986, section 2.10.

\underline{Corollary 2.9.4. Let $P,Q,G$ be commutative groups
satisfying}
\begin{equation}\tag{2.9.4.1}
 \underline{\operatorname{Hom}}_{\mathrm{pt}}(P,G)=e,
 \qquad
 \underline{\operatorname{Hom}}_{\mathrm{pt}}(Q,G)=e
\end{equation}
\underline{(cf.\ 1.3.7). Then the hypothesis of Proposition 2.9.3 is
satisfied, and conclusions (a) and (b) of that proposition hold.
Moreover, if $E$ is a torsor under $G_{P\times Q}$ birigidified
relative to $(e_P,e_Q)$, it admits a biextension structure of $(P,Q)$
by $G$ compatible with its given structures if and only if the two
morphisms of pointed sheaves which it defines,}
\begin{equation}\tag{2.9.4.2}
 Q\longrightarrow R^1p_*(G_P),
 \qquad
 P\longrightarrow R^1q_*(G_Q),
\end{equation}
\underline{where $p:P\longrightarrow e$ and $q:Q\longrightarrow e$
are the structural morphisms, are morphisms of groups.}

The first assertion was proved in 1.3.7. For the second, necessity is
clear. To prove sufficiency, apply Proposition 2.9.3(b) and 1.3.5(b).
Multiplicativity of the first morphism in (2.9.4.2) says that the
rigidified torsor $E$ under $G_P$ satisfies the condition of 1.3.5(b)
\underline{locally over $P$}. As observed in 1.3.7, that condition is
then true globally over $P$, so $E$ admits an extension structure of
$Q_P$ by $G_P$ compatible with its rigidified torsor structure.
Symmetrically, $E$ admits an extension structure of $P_Q$ by $G_Q$
compatible with its rigidified torsor structure under $G_Q$. The
conclusion follows from Proposition 2.9.3(b).

\underline{Example 2.9.5.} Let $S$ be a scheme, let $P$ and $Q$ be
abelian schemes over $S$, and let $G=\underline G_{m,S}$ be the
multiplicative group scheme over $S$. We noted in 1.3.8 that
(2.9.4.1) then holds. A birigidified torsor $E$ over $P\times Q$ is
equivalently a birigidified invertible sheaf over $P\times Q$. Recall
that a \underline{divisorial correspondence on $P,Q$} is the
isomorphism class of such a birigidified invertible sheaf (see, for
example, [1]). In the fppf topology, the morphisms (2.9.4.2) are
\begin{equation}\tag{2.9.5.1}
 Q\longrightarrow\underline{\operatorname{Pic}}_{P/S},
 \qquad
 P\longrightarrow\underline{\operatorname{Pic}}_{Q/S},
\end{equation}
defined by the divisorial correspondence. The theorem of the square
(\emph{loc.\ cit.}) says that these morphisms necessarily respect the
group structures. Thus 2.9.4 shows that
\underline{the category of biextensions of $(P,Q)$ by
$G=\underline G_{m,S}$ is equivalent to the category of birigidified
invertible sheaves on $P\times Q$}; it is therefore a rigid category,
and its set of isomorphism classes identifies with the set of
divisorial correspondences on $(P,Q)$. This identification plainly
respects the natural group laws on the two sets, the first being that
of (2.5.2.2).

\underline{Remark 2.9.6.} Additivity of the first morphism in
(2.9.4.2) is equivalent to factorization of the second through the
subsheaf $\underline{\operatorname{Ext}}^1(Q,G)$; note that, by 1.3.5,
the canonical morphism
\[
 \underline{\operatorname{Ext}}^1(Q,G)
 \longrightarrow R^1q_*(G_Q)
\]
is injective. Likewise, additivity of the second morphism in
(2.9.4.2) means that the first factors through
$\underline{\operatorname{Ext}}^1(P,G)$. This gives group
isomorphisms
\begin{equation}\tag{2.9.6.1}
 \operatorname{Biext}^{1}(P,Q;G)
 \simeq
 \operatorname{Hom}\bigl(Q,\underline{\operatorname{Ext}}^1(P,G)\bigr)
 \simeq
 \operatorname{Hom}\bigl(P,\underline{\operatorname{Ext}}^1(Q,G)\bigr).
\end{equation}
In the situation of 2.9.5 these reduce to the familiar isomorphisms
\begin{equation}\tag{2.9.6.2}
 \operatorname{Corr}(P,Q)
 \simeq \operatorname{Hom}(P,Q')
 \simeq \operatorname{Hom}(Q,P'),
\end{equation}
where
\[
 P'=\underline{\operatorname{Ext}}^1(P,\underline G_m),
 \qquad
 Q'=\underline{\operatorname{Ext}}^1(Q,\underline G_m)
\]
are the dual abelian schemes of $P$ and $Q$, identified respectively
with $\underline{\operatorname{Pic}}^0_{P/S}$ and
$\underline{\operatorname{Pic}}^0_{Q/S}$. In both formulas,
$\operatorname{Hom}$ denotes the group of \underline{group
homomorphisms}.

